\documentclass[a4paper]{article}

\usepackage{graphicx}
\usepackage{amsmath,amssymb}
\usepackage{minted}
\usepackage{multirow}
\usepackage{float}
\usepackage{todonotes}
\usepackage{forest}
\usepackage{xstring}
\usepackage{xcolor}
\usepackage[most]{tcolorbox}
\usepackage{xparse}
\usepackage{natbib}

\usetikzlibrary{graphs,graphdrawing}
\usegdlibrary{layered}

% for external graphviz
\input{/home/donkarlo/Dropbox/repo/nd_ai_project/src/nd_ai/knoweldge_representation/language/formal/computer/programming/tex/latex/code_clips/external_graphviz.tex}

% for tags
\input{/home/donkarlo/Dropbox/repo/nd_ai_project/src/nd_ai/knoweldge_representation/language/formal/computer/programming/tex/latex/parts/control_sequence/kind/semtag/semtag.tex}

% for links
\input{/home/donkarlo/Dropbox/repo/nd_ai_project/src/nd_ai/knoweldge_representation/language/formal/computer/programming/tex/latex/code_clips/seehere.tex}

\title{Remebering event specific}
\date{}

\begin{document}
    \maketitle
    \cite{conway-2000-the-construction-of-autobiographical-memories-in-the-self-memory-system} says the process of remembering an event specific knowledge is as follows by order (This paper is the base paper for remebering event specific knowledge)
    \begin{enumerate}
        \item \textbf{WillingToRemember}\\
        The system intentionally initiates an attempt to recall a past experience.
        
        \item \textbf{RetrievalModeEntering}\\
        The mind enters a retrieval state in which memory construction becomes the current active task.
        
        \item \textbf{WorkingSelfGoalsActivating}\\
        Current self-related goals become active and constrain what kind of memory can be searched for and accepted.
        
        \item \textbf{CueElaborating}\\
        The initial cue is expanded and refined to make it more effective for accessing autobiographical knowledge.
        
        \item \textbf{VerificationCriteriaSetting}\\
        The system sets criteria for what will count as an acceptable retrieved memory.
        
        \item \textbf{RetrievalModelBuilding}\\
        A retrieval model is formed that combines cues, goals, and criteria to guide the search process.
        
        \item \textbf{KnowledgeBaseSearching}\\
        The autobiographical knowledge base is searched through spreading activation and controlled access.
        
        \item \textbf{GeneralEventAccessing}\\
        General event knowledge is typically accessed first, providing an initial structured entry point into memory.
        
        \item \textbf{StableMemoryPatternForming}\\
        A stable activation pattern is formed across relevant autobiographical knowledge structures.
        
        \item \textbf{EventSpecificKnowledgeActivating}\\
        Event-specific knowledge is activated, making detailed information about the target event available.
        
        \item \textbf{SensoryDetailsRetrieving}\\
        Sensory-perceptual details such as imagery, sounds, feelings, or other experiential fragments enter awareness.
    \end{enumerate}
    
    
    A probabale data flow from goal to planning
    
    \begin{verbatim}
SurvivalGoal
        ↓
        PlanningNeed
        ↓
        AutobiographicalMemoryRetrieval
        ↓
        RememberingMode
        ├── ConceptualRemembering
        └── PerceptualRemembering
        ↓
        Planning
    \end{verbatim}
    
    
    \bibliographystyle{plainnat}
    \bibliography{/home/donkarlo/Dropbox/repo/nd_shared_research/bibliography/bibliography.bib}
\end{document}